\documentclass[12pt,letterpaper]{article}

% ---- Packages ----
\usepackage[utf8]{inputenc}
\usepackage[T1]{fontenc}
\usepackage{amsmath,amssymb}
\usepackage{newtxtext,newtxmath}  % Times New Roman text + math (replaces mathptmx)
\usepackage[margin=1in]{geometry}
\usepackage{setspace}
\usepackage{titlesec}
\usepackage{enumitem}
\usepackage{hyperref}
\usepackage{microtype}
\usepackage{booktabs}

% ---- Section formatting ----
\titleformat{\section}{\large\bfseries\sffamily}{\thesection.}{0.5em}{}
\titleformat{\subsection}{\normalsize\bfseries\sffamily}{\thesubsection}{0.5em}{}
\titleformat{\subsubsection}{\normalsize\bfseries}{\thesubsubsection}{0.5em}{}

% ---- Spacing ----
\setstretch{1.15}
\setlength{\parskip}{0.5em}
\setlength{\parindent}{0em}

% ---- Custom commands ----
\newcommand{\rhoa}{\rho_{\!A}}
\newcommand{\Ka}{K_{\!A}}

\begin{document}

\begin{center}
{\LARGE\bfseries Superfluid Cooling Feedback as the Origin\\
of Dark Energy and Dark Matter}\\[1em]
{\large Erik Otto}\\[0.5em]
{\normalsize\itshape Working Paper in the Dynamic Aether Mechanics Framework}
\end{center}

\vspace{1em}

\begin{abstract}
Within the Dynamic Aether Mechanics framework, a single mechanism---the
freeze-out of roton-like excitations as the cosmic Aether cools---is shown to
produce both dark energy and dark matter through a self-reinforcing feedback
loop.\ As the universe expands, the Aether temperature drops below the roton
energy gap, exponentially suppressing the normal-fluid fraction.\ This
collapse in normal fraction eliminates the bulk viscosity that had been
braking cosmic expansion, creating an irreversible ratchet: expansion causes
cooling, cooling reduces viscous resistance, reduced resistance permits faster
expansion.\ The mechanism naturally explains the onset of cosmic acceleration
$\sim$5~Gyr ago (the epoch when the Aether temperature crossed below the roton
gap), the phantom-to-quintessence crossing observed by DESI, and the near
equality of the CMB temperature and the transition temperature.\ Early-universe
dark matter effects arise through a complementary channel: the
temperature-dependent gravitational asymmetry fraction $\epsilon(T)$ produces
stronger effective gravity in cold, dense regions, mimicking the gravitational
signatures conventionally attributed to dark matter particles.\ The two dark
phenomena are thus revealed as dual consequences of the same underlying
thermodynamic transition---resolving the cosmic coincidence problem.
\end{abstract}

\vspace{1em}

\bigskip
\noindent\textit{This working paper explores a specific mechanism within the
Dynamic Aether Mechanics framework.\ It references results from Paper~1
(General Overview), Paper~3 (Gravity), and Paper~6 (Cosmology) of the main
series.\ The ideas presented here are under development and may be
incorporated into the main papers in revised form.}
\bigskip


% ============================================================
\section{Introduction}
% ============================================================

Modern cosmology faces two deeply puzzling observations: approximately 68\%
of the universe's energy budget consists of dark energy---a repulsive
influence causing the expansion of the universe to accelerate---while
approximately 27\% consists of dark matter---an attractive influence that
binds galaxies and clusters more tightly than their visible mass can
account for.\ The standard $\Lambda$CDM model treats these as entirely
separate phenomena: a cosmological constant $\Lambda$ for dark energy and
cold, collisionless particles for dark matter.

This separation raises an uncomfortable question known as the \emph{cosmic
coincidence problem}: why are the densities of dark energy and dark matter
similar in magnitude \emph{right now}?\ In $\Lambda$CDM, dark energy density
is constant while matter density dilutes as $a^{-3}$; their near-equality
today is a fleeting accident with no physical explanation.

The Dynamic Aether Mechanics framework (Papers~1--7) offers a different
perspective.\ In this framework, both dark energy and dark matter arise from
thermodynamic properties of a single medium---the Aether---specifically from
its two-fluid (superfluid + normal) structure.\ Paper~6 identifies dark energy
as the thermal pressure of hot void Aether, sustained by bulk-viscous
dissipation of the Hubble flow, and dark matter as the two-fluid first-sound
correction that produces additional centripetal acceleration through the
acoustic metric.

This working paper develops a specific mechanism---\emph{roton freeze-out}---that
connects these two phenomena through a single thermodynamic transition and
explains their simultaneous emergence.\ The central idea is that the Aether's
excitation spectrum includes a roton-like energy gap $\Delta$, and as the
cosmic temperature drops below $\Delta/k_B$, the resulting exponential
suppression of the normal fraction triggers a cascade of consequences that
manifest as both dark energy and dark matter.


% ============================================================
\section{The Cooling Feedback Mechanism}
\label{sec:cooling}
% ============================================================

\subsection{Global Cooling of the Free Aether}

Two processes have shaped the thermal state of the Aether over cosmic history
(Paper~6, Section~1):

\begin{enumerate}[leftmargin=2em]
  \item \textbf{Permanent binding} locks thermal energy into mass.\ Each
    mass-formation event (driven by the Dynamic Casimir Effect) extracts energy
    from the local Aether, cooling the matter-containing regions.

  \item \textbf{Stellar radiation} heats the Aether in voids, partially
    offsetting the cooling.\ However, the integrated stellar radiation
    ($\sim$$10^{-14}$~J/m$^3$) falls $\sim$$10^4$ times short of the observed
    dark energy density, so the dominant energy channel is bulk-viscous
    dissipation of the Hubble flow.
\end{enumerate}

The net effect over cosmic time is a gradual decrease in the mean Aether
temperature.\ The CMB temperature $T_{\mathrm{CMB}}(z) = T_0(1 + z)$, where
$T_0 = 2.725$~K, tracks this cooling: at redshift $z = 2$ the Aether was at
$\sim$8~K, and at $z = 5$ it was at $\sim$16~K.


\subsection{Roton-Like Excitations and the Energy Gap}

In superfluid helium-4, the excitation spectrum contains two branches: low-energy
phonons ($\omega = ck$) and higher-energy rotons with a characteristic energy gap
$\Delta_{\mathrm{He}}/k_B \approx 8.6$~K.\ Below this gap, roton excitations are
exponentially suppressed, and the normal fraction collapses:
%
\begin{equation}
  \frac{\rho_n}{\rhoa} \;\propto\; \exp\!\left(-\frac{\Delta}{k_B T}\right)
  \quad\text{for}\quad T \lesssim \Delta/k_B
  \label{eq:normal-fraction}
\end{equation}

If the Aether possesses analogous intermediate-wavelength excitations (an open
question identified in Paper~6, Section~1.10), the same exponential suppression
occurs as the cosmic temperature drops below the gap energy.\ The roton gap
$\Delta/k_B$ for the Aether need not equal helium's value; it is determined by
the Aether's specific interaction parameters ($m_A$, $\rhoa$, $\Ka$).

Note that the full Landau formula for helium-4 rotons includes a prefactor
$\rho_{n,\mathrm{roton}} \propto T^{-1/2}\exp(-\Delta/k_B T)$.\ The
$T^{-1/2}$ factor is specific to helium's roton dispersion relation; the
Aether's excitation spectrum may produce a different prefactor.\ The
exponential Boltzmann suppression is universal for any gapped excitation and
dominates the temperature dependence.\ Throughout this paper, the pure
exponential is used as an approximation; the prefactor introduces corrections
of order $(1+z)^{-1/2}$ (a $\sim$40\% reduction at $z = 2$) that do not
qualitatively change the results.

Paper~6 already identifies the excitation spectrum as the critical unknown that
simultaneously determines the dark matter magnitude ($\Delta c_1^2/c^2 \sim
10^{-6}$), the bulk viscosity ($\zeta_n \sim 10^6$~Pa\,s), the void temperature
ratio ($T_{\mathrm{void}}/T_{\mathrm{wall}} \approx 2.5$), and the dark energy
equation of state $w(z)$.\ This paper proposes that the roton gap $\Delta$ is the
single parameter that governs all four.


\subsection{The Self-Reinforcing Ratchet}

The cooling feedback operates through a self-reinforcing cycle:

\begin{enumerate}[leftmargin=2em]
  \item \textbf{Expansion does adiabatic work.}\ The Hubble flow cools the
    normal component of the Aether (the superfluid component, carrying zero
    entropy, cannot be cooled further).

  \item \textbf{Cooling suppresses roton excitations.}\ As $T$ drops further
    below $\Delta/k_B$, roton modes freeze out exponentially, converting normal
    fluid to superfluid.

  \item \textbf{The feedback becomes self-sustaining.}\ In principle, heating
    can reconvert superfluid to normal fluid (as in helium-4).\ However, the
    viscous heating that could replenish the normal fraction is itself
    \emph{proportional to the normal fraction} ($\zeta_n \propto \rho_n$).\
    Once $\rho_n$ drops below a critical threshold, the heating rate becomes
    too weak to offset the cooling from expansion---the feedback loop becomes
    supercritical.\ From this point, cooling outpaces heating, the normal
    fraction continues to fall, and the process is dynamically irreversible:
    a self-reinforcing instability rather than a thermodynamic one-way door.

  \item \textbf{Viscosity collapses.}\ Bulk viscosity $\zeta_n$ is carried
    exclusively by the normal component.\ As the normal fraction drops, so does
    $\zeta_n$, reducing the viscous braking that had been decelerating the
    expansion.

  \item \textbf{Expansion accelerates.}\ With less viscous resistance, the
    existing expansion---driven by the thermal pressure differential between
    voids and matter regions---encounters less opposition.\ The expansion rate
    increases.

  \item \textbf{Return to step 1.}\ Faster expansion produces more adiabatic
    cooling, driving further roton freeze-out.
\end{enumerate}

This is a positive feedback loop that becomes self-sustaining once the normal
fraction drops below the critical threshold (step~3).\ Below this threshold,
the ratchet only tightens: the system cannot return to the high-viscosity,
braking-dominated regime because the mechanism that would restore it (viscous
heating) has been weakened by the very transition it would need to reverse.\
The result is an accelerating expansion that we observe as dark energy.


% ============================================================
\section{Dark Energy from Roton Freeze-Out}
\label{sec:dark-energy}
% ============================================================

\subsection{Viscous Heating Rate}

Paper~6 derives the bulk-viscous dissipation rate from the Hubble flow:
%
\begin{equation}
  \dot\varepsilon \;=\; 9\,\zeta_n\,H^2
  \label{eq:dissipation}
\end{equation}

In the roton freeze-out model, the bulk viscosity evolves with redshift through
the normal fraction:
%
\begin{equation}
  \zeta_n(z) \;=\; \zeta_0 \, f_n(z)
  \label{eq:zeta-evolution}
\end{equation}
%
where $\zeta_0 \sim 10^6$~Pa\,s is the present-day value and the relative
normal fraction is:
%
\begin{equation}
  f_n(z) \;\equiv\; \frac{\rho_n(z)}{\rho_n(0)}
    \;=\; \exp\!\left[\frac{\Delta}{k_B T_0}\,\frac{z}{1+z}\right]
  \label{eq:fn}
\end{equation}

This follows from Eq.~\ref{eq:normal-fraction} with $T(z) = T_0(1+z)$:
%
\begin{equation}
  f_n(z) = \frac{\exp(-\Delta/[k_B T_0(1+z)])}{\exp(-\Delta/[k_B T_0])}
    = \exp\!\left[\frac{\Delta}{k_B T_0}\left(1 - \frac{1}{1+z}\right)\right]
    = \exp\!\left[\frac{\alpha\,z}{1+z}\right]
\end{equation}
%
where $\alpha \equiv \Delta/(k_B T_0)$.


\subsection{Normal Fraction and Heating Rate Tables}

For $\Delta/k_B = 6$~K (constrained below by the DESI crossing redshift),
$\alpha = 6/2.725 = 2.20$.\ The resulting evolution:

\begin{table}[h]
\centering
\begin{tabular}{@{}lccccc@{}}
\toprule
Redshift $z$ & $T_{\mathrm{CMB}}$ (K) & $f_n(z)$ &
  $(H/H_0)^2$ & $S(z)/S(0)$ \\
\midrule
0   & 2.73 & 1.00 & 1.00  & 1.0 \\
0.3 & 3.54 & 1.66 & 1.36 & 2.3 \\
0.5 & 4.09 & 2.08 & 1.71 & 3.6 \\
0.7 & 4.63 & 2.48 & 2.17 & 5.4 \\
1.0 & 5.45 & 3.01 & 3.10 & 9.3 \\
2.0 & 8.18 & 4.34 & 8.80 & 38.2 \\
5.0 & 16.4 & 6.26 & 65.5 & 410 \\
\bottomrule
\end{tabular}
\caption{Evolution of the normal fraction ratio $f_n(z)$, the squared Hubble
  parameter, and the total viscous heating rate $S(z) = 9\zeta_n(z)H(z)^2$
  relative to today, for a roton gap $\Delta/k_B = 6$~K.\ The Hubble parameter
  uses $H(z) = H_0\sqrt{0.3(1+z)^3 + 0.7}$.}
\label{tab:heating}
\end{table}

The viscous heating rate at $z = 1$ was approximately 9$\times$ stronger than
today, and at $z = 2$ it was approximately 38$\times$ stronger.\ This dramatic
change is driven primarily by the exponential roton factor at moderate redshifts
and by the $H^2$ factor at higher redshifts.


\subsection{Phantom--Quintessence Crossing}

The effective dark energy equation of state parameter $w(z)$ is defined by how
the dark energy density $\rho_{\mathrm{DE}}$ evolves:
%
\begin{equation}
  w(z) \;=\; -1 \;-\; \frac{1}{3H}\,\frac{d\ln\rho_{\mathrm{DE}}}{dt}
  \label{eq:w-definition}
\end{equation}

The dark energy density evolves through two competing effects:

\begin{itemize}[leftmargin=2em]
  \item \textbf{Viscous heating} (adds energy):
    $\dot\varepsilon_{\mathrm{heat}} = 9\,\zeta_n(z)\,H^2$
  \item \textbf{Expansion dilution} (removes energy):
    $\dot\varepsilon_{\mathrm{dilute}} \propto H\,P_{\mathrm{th}}$
    (see Section~\ref{sec:wz} for discussion of the dilution rate)
\end{itemize}

The balance determines the sign of $d\rho_{\mathrm{DE}}/dt$ and hence
whether $w$ is above or below $-1$:

\begin{itemize}[leftmargin=2em]
  \item When heating dominates ($S > D$): $\rho_{\mathrm{DE}}$ is increasing
    $\Rightarrow$ $w < -1$ \textbf{(phantom)}
  \item When dilution dominates ($D > S$): $\rho_{\mathrm{DE}}$ is decreasing
    $\Rightarrow$ $w > -1$ \textbf{(quintessence)}
  \item At the crossover ($S = D$): $w = -1$ exactly
\end{itemize}

At high redshift, the viscous heating rate was far stronger than today
(Table~\ref{tab:heating}), so heating dominated---placing the early dark energy
firmly in the phantom regime ($w < -1$).\ As the Aether cooled through the roton
gap, the normal fraction collapsed exponentially, viscous heating dropped, and
dilution began to dominate---transitioning to the quintessence regime ($w > -1$).

This is precisely the phantom-to-quintessence crossing observed by DESI at
2.8--4.2$\sigma$ significance \cite{desi2025}.


% ============================================================
\section{The \texorpdfstring{$w(z)$}{w(z)} Equation of State}
\label{sec:wz}
% ============================================================

\subsection{Evolution of the Dark Energy Density}

The thermal pressure $P_{\mathrm{th}}$ of the void Aether evolves through two
competing effects:
%
\begin{equation}
  \frac{dP_{\mathrm{th}}}{dt} \;=\;
    \underbrace{9\,\zeta_n(z)\,H^2}_{\text{viscous heating}}
    \;-\; \underbrace{\beta\,H\,P_{\mathrm{th}}}_{\text{expansion work}}
  \label{eq:P-evolution}
\end{equation}
%
where $\beta$ parameterizes the rate at which the thermal pressure does
work on the expansion.\ In a standard expanding fluid, $\beta = 3$
(matter-like dilution $\propto a^{-3}$).\ However, the Aether density is
held essentially constant ($\delta\rhoa/\rhoa \sim 10^{-4}$) by the
enormous bulk modulus $\Ka \sim 10^{28}$~J/m$^3$, so the thermal energy
density does not dilute in the usual sense.\ The effective $\beta$ depends
on how efficiently the thermal pressure converts to kinetic energy of expansion
and is an open calculation within the Aether framework.

\paragraph{Why the standard $w(z)$ mapping is non-trivial.}
In the Aether model, dark energy is not a homogeneous negative-pressure fluid
as in $\Lambda$CDM.\ It is the \emph{thermal pressure gradient} between hot
voids and cold matter-containing regions that drives galaxy recession---a real
force, not a geometric effect.\ The equation of state parameter $w(z)$ is an
\emph{effective} quantity extracted when observers fit the expansion history
$H(z)$ to the $\Lambda$CDM framework.\ Deriving $w(z)$ precisely requires
computing $H(z)$ from the Aether's pressure-gradient dynamics and then
fitting it to the standard parameterization---a calculation that depends on
$\beta$ and the full thermal history.

\paragraph{Robust qualitative prediction.}
Despite the mapping complexity, the \emph{direction} of the phantom--quintessence
crossing is robust and does not depend on $\beta$.\ The viscous heating rate
$S(z) = 9\,\zeta_n(z)\,H^2$ was dramatically stronger in the past
(Table~\ref{tab:heating}).\ Regardless of the dilution rate, this means:

\begin{itemize}[leftmargin=2em]
  \item At high $z$: strong viscous heating was building up the effective dark
    energy density faster than dilution could reduce it $\Rightarrow$ when fit to
    $\Lambda$CDM, this appears as $w < -1$ \textbf{(phantom)}.
  \item At low $z$: roton freeze-out has collapsed the heating rate, so whatever
    dilution exists now dominates $\Rightarrow$ this appears as $w > -1$
    \textbf{(quintessence)}.
  \item The crossover from one regime to the other is the phantom--quintessence
    crossing observed by DESI.
\end{itemize}


\subsection{Constraining the Roton Gap from DESI}

The phantom--quintessence crossing occurs at the redshift $z_*$ where the
viscous heating rate equals the expansion dilution rate
(Eq.~\ref{eq:P-evolution} with $dP_{\mathrm{th}}/dt = 0$).\ The crossing
redshift depends on $\alpha = \Delta/(k_B T_0)$ through
$\zeta_n(z) = \zeta_0\,\exp[\alpha z/(1+z)]$.\ Since $\zeta_n$ grows
exponentially with $z$ while the dilution rate grows more slowly, the
crossing moves to higher $z$ for smaller $\alpha$ (smaller roton gap---slower
exponential growth, so the balance point shifts outward).

\begin{table}[h]
\centering
\begin{tabular}{@{}lcc@{}}
\toprule
$\Delta/k_B$ (K) & $\alpha$ & Estimated $z_*$ \\
\midrule
4  & 1.47 & $\sim$1.5 \\
5  & 1.83 & $\sim$1.0--1.2 \\
6  & 2.20 & $\sim$0.7--1.0 \\
8  & 2.94 & $\sim$0.3--0.5 \\
12 & 4.40 & $\sim$0.1--0.2 \\
\bottomrule
\end{tabular}
\caption{Estimated phantom--quintessence crossing redshift $z_*$ as a function
  of the roton gap $\Delta/k_B$.\ DESI observations indicate $z_* \gtrsim 1$,
  constraining $\Delta/k_B \approx 5$--$6$~K.}
\label{tab:crossing}
\end{table}

DESI's observed crossing at $z \gtrsim 1$ constrains the roton gap to
$\Delta/k_B \approx 5$--$6$~K.\ This is modestly below the helium-4 roton gap
($8.6$~K), consistent with the fact that the Aether has fundamentally different
microscopic parameters ($m_A \sim 700$~MeV/$c^2$ vs.\ $m_{\mathrm{He}} =
4$~amu; $\rhoa \sim 10^{12}$~kg/m$^3$ vs.\ $\rho_{\mathrm{He}} \sim 145$~kg/m$^3$).


\subsection{Qualitative Shape of \texorpdfstring{$w(z)$}{w(z)}}

The model predicts a distinctive $w(z)$ shape:

\begin{itemize}[leftmargin=2em]
  \item At $z \gg z_*$: viscous heating strongly dominates $\Rightarrow$
    $w \ll -1$ (deep phantom)
  \item Near $z \approx z_*$: heating and dilution balance $\Rightarrow$
    $w$ crosses $-1$
  \item At $z = 0$: dilution dominates $\Rightarrow$ $w > -1$ (quintessence)
  \item Far future ($z < 0$): roton freeze-out essentially complete,
    $\zeta_n \to 0$ $\Rightarrow$ $w$ approaches a value determined by the
    dilution rate $\beta$ (no viscous heating to offset dilution)
\end{itemize}

Crucially, $w(z)$ follows a \emph{phase-transition curve} (exponential in
$\Delta/T$), not the smooth polynomial assumed by the CPL parameterization
$w(a) = w_0 + w_a(1-a)$.\ Future high-precision $w(z)$ measurements could
distinguish between the exponential form predicted here and polynomial models.


\subsection{The CMB Temperature Coincidence}

The current CMB temperature ($T_0 = 2.725$~K) lies within a factor of two of
the constrained roton gap ($\Delta/k_B \approx 5$--$6$~K).\ This is not a
coincidence---it is a \emph{consequence} of the fact that we observe the universe
during the transition epoch.\ Dark energy became observationally significant when
the normal fraction dropped enough for the feedback ratchet to dominate.\ We
detect it now precisely because $T_0$ is in the range where roton freeze-out
is actively occurring.

The ``coincidence problem'' of $\Lambda$CDM---why is dark energy dominant
\emph{now}?---becomes a phase-transition timing question with a physical answer:
the Aether crossed $\Delta/k_B$ in the recent cosmological past, and the
resulting ratchet is still tightening.


% ============================================================
\section{Dark Matter from Varying \texorpdfstring{$G$}{G}}
\label{sec:dark-matter}
% ============================================================

\subsection{Temperature Dependence of the Gravitational Constant}

Paper~3 derives the gravitational constant as $G \propto \epsilon\,\rhoa\,\psi^2$,
where $\epsilon$ is the \emph{asymmetry fraction}---the fraction of each transient
bind--unbind cycle's energy that escapes as the gravitational wave.\ Since the
Aether density $\rhoa$ is essentially constant ($\delta\rhoa/\rhoa \sim 10^{-4}$),
the primary channel for $G$ variation is through $\epsilon(T)$:

\begin{quote}
Higher temperature provides thermal energy that resists binding, potentially
reducing the fraction of each cycle's energy that escapes as the gravitational
wave.\ Cold environments (near matter) could have higher $\epsilon$ and therefore
stronger gravitational output, while hot environments (cosmic voids) could have
lower $\epsilon$.\ (Paper~3, Section~1.3)
\end{quote}

This means:
%
\begin{equation}
  G(T) \;=\; G_0 \;\frac{\epsilon(T)}{\epsilon(T_0)}
  \label{eq:G-varying}
\end{equation}
%
where $G_0$ is the locally measured value (in cold, matter-containing regions).


\subsection{Magnitude of the Variation}

For structure formation, the dark matter mass-to-visible-mass ratio is
approximately 5--6:1.\ If varying $G$ is responsible, we need:
%
\begin{equation}
  \frac{G_{\mathrm{dense}}}{G_{\mathrm{void}}}
    \;=\; \frac{\epsilon(T_{\mathrm{matter}})}{\epsilon(T_{\mathrm{void}})}
    \;\approx\; 6
  \label{eq:G-ratio}
\end{equation}

With the void-to-matter temperature ratio $T_{\mathrm{void}}/T_{\mathrm{matter}}
\approx 2.5$ (Paper~6, Eq.~18), this requires:
%
\begin{equation}
  \epsilon(T) \;\propto\; T^{-n}, \qquad
  2.5^n = 6 \;\Rightarrow\; n \approx 1.95
  \label{eq:epsilon-power}
\end{equation}

A near-quadratic temperature dependence $\epsilon \propto T^{-2}$ is physically
motivated: each transient binding event involves a \emph{pair} of particles
(particle--antiparticle pair from the Dynamic Casimir Effect), and the thermal
disruption of each member's binding efficiency contributes independently,
giving a product of two temperature-dependent factors.


\subsection{Consistency with Local \texorpdfstring{$G$}{G} Measurements}

The variation $G \propto T^{-2}$ does not conflict with the exquisite precision
of local $G$ measurements.\ All terrestrial and solar-system measurements of $G$
occur in the same thermal environment---cold, matter-containing regions where
$T \approx T_{\mathrm{matter}}$ is essentially uniform.\ The measured value is
$G(T_{\mathrm{matter}})$, which is constant across all local experiments.

The \emph{spatial} variation of $G$ between galaxy cores and cosmic voids
($\Delta G/G \sim 5$--$6$) is not directly measured, because:

\begin{itemize}[leftmargin=2em]
  \item Gravitational measurements (rotation curves, lensing) are made
    \emph{near matter}, where $G = G_0$.
  \item What is conventionally attributed to ``dark matter mass'' is the
    \emph{difference} between the gravitational effect produced by
    $G_{\mathrm{dense}}$ and what would be expected from $G_{\mathrm{void}}$
    applied to the same visible mass.
  \item No experiment directly compares $G$ at two locations with different
    Aether temperatures.
\end{itemize}


\subsection{Constraints from the CMB}

At the CMB epoch ($z \sim 1100$, $T \approx 3000$~K), the Aether temperature
was nearly uniform: $\delta T/T \sim 10^{-5}$.\ The corresponding $G$ variation:
%
\begin{equation}
  \frac{\delta G}{G} \;=\; n\,\frac{\delta T}{T}
    \;\approx\; 2 \times 10^{-5}
  \label{eq:deltaG-CMB}
\end{equation}

This is a 0.002\% variation---well within current CMB constraints on $G$
variability, which are at the percent level.\ The effect would produce
subtle modifications to the CMB high-$\ell$ power spectrum (enhanced Silk
damping in slightly warmer regions), potentially detectable by next-generation
polarization experiments (CMB-S4, LiteBIRD).


% ============================================================
\section{Mechanism Handoff: Early to Late Universe}
\label{sec:handoff}
% ============================================================

The superfluid cooling model predicts that the \emph{mechanism} underlying
``dark matter'' effects changes smoothly over cosmic time, even though the
observational signatures remain broadly similar.

\begin{table}[h]
\centering
\begin{tabular}{@{}lccll@{}}
\toprule
Epoch & $z$ & $T_{\mathrm{CMB}}$ (K) &
  Primary DM mechanism & Primary DE mechanism \\
\midrule
Recombination & $\sim$1100 & $\sim$3000 & First-sound correction & Viscous braking (strong) \\
Structure formation & 2--20 & 8--55 & First-sound + varying $G$ & Braking (strong) \\
Transition & 0.5--2 & 4--8  & Both strengthening & Braking weakening \\
Present & 0 & 2.7 & First-sound (peak) & Ratchet engaged \\
Future & $<$0 & $<$2.7 & Weakening & Ratchet complete \\
\bottomrule
\end{tabular}
\caption{The dominant dark matter and dark energy mechanisms at different
  cosmic epochs.\ The first-sound correction (Paper~6) operates at all
  epochs.\ Varying $G$ becomes a significant \emph{supplementary} channel
  once structures form and temperature contrasts develop ($z \lesssim 20$).}
\label{tab:handoff}
\end{table}

\paragraph{Recombination ($z \sim 1100$, $T \sim 3000$~K).}
At the CMB epoch, the universe is nearly homogeneous ($\delta T/T \sim 10^{-5}$),
so varying $G$ is negligible ($\delta G/G \sim 2 \times 10^{-5}$).\ Dark matter
effects at this epoch---reflected in the CMB acoustic peak heights---arise from
the first-sound correction (Paper~6), which operates whenever a temperature
gradient exists between denser and less dense regions.\ The derivation of the
CMB power spectrum from first-sound dynamics remains an open calculation
(Paper~6, Section~1.10).

\paragraph{Structure formation ($z \sim 2$--$20$, $T \sim 8$--$55$~K).}
As proto-structures collapse, they cool the local Aether through binding,
creating growing temperature contrasts with the surrounding medium.\ This
bootstrapping process progressively enhances $G$ in dense regions: slight
cooling $\to$ slightly stronger $G$ $\to$ faster collapse $\to$ more cooling.\
By $z \sim 2$--$3$, significant structures exist with $T_{\mathrm{void}}/
T_{\mathrm{matter}} \approx 2.5$, and the varying-$G$ channel contributes
substantially alongside the first-sound correction.

\paragraph{Transition ($z \sim 0.5$--$2$, $T \sim 4$--$8$~K).}
The Aether is cooling through the roton gap.\ The normal fraction is dropping
exponentially.\ Both varying $G$ (still significant where temperature contrasts
persist) and the first-sound correction (strengthening as the superfluid
fraction increases near matter while voids retain more normal component)
contribute to dark matter effects.\ Dark energy is ``turning on'' as viscous
braking collapses.

\paragraph{Present ($z \approx 0$, $T \approx 2.7$~K).}
The normal fraction is low globally.\ Near matter (cold regions), the Aether
is nearly pure superfluid.\ In voids (warmer), some normal component persists.\
The contrast in $c_1$ between these regions is at its peak, making the
first-sound correction the dominant dark matter mechanism for galaxy rotation
curves and gravitational lensing.\ Dark energy is fully engaged---the ratchet
is tightening.

\paragraph{Far future ($z < 0$, $T < 2.7$~K).}
Eventually, even the voids go fully superfluid.\ The $c_1$ gradient flattens,
weakening dark matter effects.\ $G$ variation also weakens (temperatures
converge).\ Galaxies gradually become less bound.\ Expansion continues
accelerating (the ratchet is self-sustaining), heading toward a fully superfluid
cosmos.


% ============================================================
\section{Observational Consistency Checks}
\label{sec:checks}
% ============================================================

\subsection{DESI Dark Energy Equation of State}

\textbf{Observation:} $w_0 > -1$ and $w_a < 0$ at 2.5--3.9$\sigma$; phantom-to-quintessence crossing at $z \gtrsim 1$ \cite{desi2025}.

\textbf{Model prediction:} At $z = 0$, dilution dominates over viscous heating
($w > -1$, quintessence).\ Going back in time, the exponentially larger
$\zeta_n$ makes heating dominate ($w < -1$, phantom).\ With $\Delta/k_B \approx
5$--$6$~K, the crossing occurs at $z \approx 0.7$--$1.0$.

\textbf{Status:} \textbf{Consistent.}


\subsection{Bulk Viscosity}

\textbf{Observation:} Bulk viscous cosmology models require $\zeta_0 \sim
10^6$~Pa\,s \cite{viscous2025}.

\textbf{Model prediction:} $\zeta_0 \sim 10^6$~Pa\,s (Paper~6, Eq.~2).

\textbf{Status:} \textbf{Exact match.}


\subsection{JWST High-Redshift Galaxies}

\textbf{Observation:} Unexpectedly massive, mature galaxies at $z > 6$.\ No
clear evidence for modified dark matter properties, though results are
consistent with both CDM and WDM.

\textbf{Model prediction:} At $z > 6$, the first-sound correction provides
baseline dark matter effects, while varying $G$ (enhanced in cold, dense
regions via the bootstrapping process of Section~\ref{sec:handoff})
provides additional gravitational enhancement.\ Together, these produce
\emph{faster} structure formation than standard gravity alone.\ Massive
galaxies forming early is a natural consequence.\ The gravitational signatures
are indistinguishable from CDM because both ``more mass'' and ``stronger $G$''
produce the same observables (deeper potential wells, stronger lensing).

\textbf{Status:} \textbf{Consistent; potentially explanatory.}


\subsection{CMB Constraints on Varying \texorpdfstring{$G$}{G}}

\textbf{Observation:} High-$\ell$ CMB power spectrum constrains $G$ variations
at the percent level.

\textbf{Model prediction:} At the CMB epoch ($z \sim 1100$, $T \approx 3000$~K),
temperature variations are $\delta T/T \sim 10^{-5}$.\ The corresponding
$\delta G/G \approx 2 \times 10^{-5}$---three orders of magnitude below
current constraints.

\textbf{Status:} \textbf{No conflict.}


\subsection{Dark Matter Halo Evolution}

\textbf{Observation:} Halo concentrations at $z \sim 3$ are roughly constant
across mass ranges ($c \sim 3.5$--$4$), with weak mass dependence that
strengthens at lower $z$.

\textbf{Model prediction:} At $z \sim 3$, the dark matter effect includes both
the first-sound correction and a growing contribution from $G$ variation.\
The $G$-variation component is temperature-dependent (not mass-dependent),
producing weak mass dependence.\ At $z \sim 0$, the first-sound correction
dominates and depends on the galaxy's luminosity (which sets the temperature
profile), introducing stronger mass dependence.\ The transition from
weak to strong mass dependence is qualitatively consistent with the observed
trend.

\textbf{Status:} \textbf{Consistent.}


\subsection{Cosmic Coincidence}

\textbf{Problem:} Why are $\Omega_{\mathrm{DE}}$ and $\Omega_{\mathrm{DM}}$
similar in magnitude today?

\textbf{Resolution:} Both phenomena are consequences of the same thermodynamic
transition.\ Dark energy ``turns on'' when the normal fraction drops enough for
the ratchet to engage.\ The first-sound dark matter correction reaches its
maximum strength during the same transition (maximum contrast between matter
and void superfluid fractions).\ Their co-emergence from a single event naturally
produces similar magnitudes without fine-tuning.

\textbf{Status:} \textbf{Resolved in principle.}


% ============================================================
\section{Predictions}
\label{sec:predictions}
% ============================================================

\begin{enumerate}[leftmargin=2em]
  \item \textbf{Roton gap $\Delta/k_B \approx 5$--$6$~K.}\ This is a specific,
    testable prediction.\ If the Aether's excitation spectrum can be derived
    from first principles (from $m_A$, $\rhoa$, $\Ka$), the resulting roton gap
    must fall in this range.\ This same gap must simultaneously produce the
    correct dark matter magnitude, bulk viscosity, and void temperature ratio.

  \item \textbf{$w(z)$ follows a phase-transition curve.}\ The equation of state
    is not a smooth polynomial but contains the exponential factor
    $\exp[\Delta/(k_B T_0) \cdot z/(1+z)]$.\ High-precision $w(z)$ measurements
    from future DESI data releases, Euclid, and the Vera Rubin Observatory
    could distinguish this shape from the CPL parameterization.

  \item \textbf{Dark matter signatures evolve with redshift.}\ Specifically,
    the ratio of ``dark'' to visible mass should show subtle redshift dependence:
    at $z > 2$ (where $G$ variation contributes significantly alongside the
    first-sound correction), dark matter effects may exhibit different scaling
    relations than at $z < 1$ (where the first-sound correction alone
    dominates).\ JWST rotation curve measurements at $z \sim 2$--$3$
    could test this.

  \item \textbf{Void regions are ``further along'' the transition.}\ Voids
    with lower temperatures should show stronger expansion rates and different
    dark matter characteristics than regions near large structures.\ This
    prediction is testable with upcoming void catalogs from DESI and Euclid.

  \item \textbf{Late-time viscosity decrease.}\ The bulk viscosity of the cosmic
    medium should be measurably lower at $z = 0$ than at $z = 1$.\ This could
    manifest as redshift-dependent damping of baryon acoustic oscillations or
    gravitational wave propagation.

  \item \textbf{Massive early galaxies from enhanced $G$.}\ The model predicts
    that structure formation was accelerated by $G$ enhancement in cold, dense
    regions.\ This provides a natural explanation for the JWST observations of
    unexpectedly massive galaxies at $z > 6$, without requiring exotic dark
    matter properties or modified initial conditions.
\end{enumerate}


% ============================================================
\section{Open Questions}
\label{sec:open}
% ============================================================

\begin{enumerate}[leftmargin=2em]
  \item \textbf{Functional form of $\epsilon(T)$.}\ The assumption $\epsilon
    \propto T^{-2}$ is motivated by the pair nature of transient binding, but a
    first-principles derivation from the Dynamic Casimir Effect threshold
    conditions is needed.\ Alternative forms (exponential suppression,
    logarithmic dependence) would change the magnitude and scaling of the
    varying-$G$ dark matter effect.

  \item \textbf{Derivation of the roton gap from Aether parameters.}\ Can
    $\Delta/k_B \approx 5$--$6$~K be derived from the known Aether parameters
    ($m_A \approx 500$--$850$~MeV/$c^2$, $\rhoa \approx 10^{11}$--$10^{12}$~kg/m$^3$,
    $\Ka \approx 10^{28}$--$10^{29}$~J/m$^3$)?\ In helium-4, the roton gap
    arises from short-range correlations in the liquid structure; the analogous
    calculation for the Aether requires knowledge of the inter-quasiparticle
    potential.

  \item \textbf{Numerical $w(z)$ solution.}\ Computing the effective $w(z)$
    requires two steps: (a)~determine the dilution parameter $\beta$ from
    the Aether's constant-density thermodynamics, and (b)~solve
    Eq.~\ref{eq:P-evolution} numerically with the evolving $\zeta_n(z)$ and
    $H(z)$ to produce $H(z)$, then fit the result to the CPL parameterization
    to extract $w_0$ and $w_a$.\ Overlaying this with DESI data would provide
    a quantitative test and refine the constraint on $\Delta/k_B$.

  \item \textbf{CMB power spectrum.}\ How does the varying-$G$ mechanism at
    $z \sim 1100$ affect the acoustic peaks?\ The $\delta G/G \sim 10^{-5}$
    variation is small but could produce detectable modifications to the
    damping tail.

  \item \textbf{Rotation curves at different redshifts.}\ The handoff from
    $G$-variation to first-sound correction should produce measurable differences
    in rotation curve shapes at $z \sim 2$--$3$ versus $z \sim 0$.\ Quantitative
    predictions require the three-dimensional temperature profile $T(r,z)$ at
    different cosmic epochs.

  \item \textbf{Reconciliation with existing Paper~6 framework.}\ The thermal
    pressure mechanism (Paper~6) and the roton freeze-out mechanism (this paper)
    are complementary: thermal pressure provides the driving force for expansion,
    while roton freeze-out explains why the expansion \emph{accelerates}.\ A
    unified treatment should demonstrate that both effects arise naturally from
    the same two-fluid dynamics.
\end{enumerate}


% ============================================================
% REFERENCES
% ============================================================
\begin{thebibliography}{99}

\bibitem{desi2025}
DESI Collaboration,
``DESI DR2 Results: Measurement of the Baryon Acoustic Oscillation
scale and Hubble parameter across 11 billion years,''
\textit{arXiv:2503.14738} (2025).

\bibitem{viscous2025}
A.~Hern\'{a}ndez-Almada \textit{et al.},
``Bulk viscous cosmological models with a cosmological constant:
Observational constraints,''
\textit{arXiv:2508.11614} (2025).

\bibitem{otto-overview}
E.~Otto,
``Unifying Theory of Dynamic Aether Mechanics: General Overview,''
this series, Paper~1.

\bibitem{otto-gravity}
E.~Otto,
``Unifying Theory of Dynamic Aether Mechanics: Gravity,''
this series, Paper~3.

\bibitem{otto-cosmology}
E.~Otto,
``Unifying Theory of Dynamic Aether Mechanics: Cosmology,''
this series, Paper~6.

\end{thebibliography}

\end{document}
